\section*{Red-team нулевого промта}

Исходный нулевой промт условно восстанавливает
\(\mathbb{RP}^3\times S^1\), но заранее фиксирует кватернионный
трёхмерный групповой сектор, \(Z_2\)-факторизацию, периодический круг и
прямое произведение.

Реальная структура \(J\) сама по себе не отождествляет антиподальные точки
базы. Модулярная группа Томиты--Такесаки параметризована
\(t\in\mathbb R\), а конечномерное KMS-состояние может иметь несоизмеримые
частоты. Например,
\[
 \rho=Z^{-1}\operatorname{diag}(1,e^{-1},e^{-\sqrt2})
\]
не допускает общего ненулевого периода.

Даже после выбора \(\mathbb{RP}^3\) и кругового волокна остаются два класса
\(U(1)\)-расслоений, поскольку
\[
 H^2(\mathbb{RP}^3;\mathbb Z)=\mathbb Z_2.
\]
Следовательно, рабочий носитель остаётся сильным условным кандидатом, но его
неизбежность требует вариационного функционала и доказательства уникального
устойчивого минимума.